\begin{abstract}
Linear relational embeddings (LREs) approximate a language
model's subject-to-object mapping with an affine
function. Their faithfulness varies enormously across
relations, including weak performance on \textsc{person\_father}
and stronger performance on \textsc{country\_capital\_city}. We investigate
the factors that contribute to this variation and
demonstrate that it can be significantly predicted from a
readout that is not aware of the underlying relationship. To
do so, we propose a novel form of the Jacobian lens, the
\emph{contrastive entity lens}. Whereas the Jacobian lens is estimated
across all tokens and concepts, the entity lens
$\mbox{EL}_t$
is trained
to predict what the model will output for an entity of
type $t$; e.g., we train ELs for countries and for persons.
After estimating an entity lens, we measure how
highly it ranks the model's answer for a relation -- without
using any relation-specific information in estimation and
prediction.  Across 28 relations of the LRE dataset on
Olmo-3-7B this relation-agnostic score correlates with LRE
faithfulness at $\rho = \headlinenumber$ ($p=\headlinep$),
using all 9,747 original rows, including demonstrations and references.
We report the same full-evaluation comparison for six additional models.  This association supports a role for information already carried by
the subject representation, while leaving the causal mechanism open.  At the fact level, generic first tokens and changes in the model's
answer across demonstration contexts complicate interpretation. We
report these limitations, the model-specific tokenization and prompt
corrections, and readout baselines under the same full evaluation.
\end{abstract}
